\begin{enumerate}

\item Se proyecta la velocidad tangencial.

\[\fitmath{\mathrm{ct}=\mathrm{Vel}\cdot \sin\left(\mathrm{ang}\right)}\]

\[\fitmath{\mathrm{ct}=(590)\cdot \sin\left((1.117010721)\right)\approx 530.2885\,\mathrm{m/s}}\]

\item El óptimo de la parábola está en U=Ctheta/2.

\[\fitmath{\mathrm{uopt}=\frac{\mathrm{Vel}\cdot \sin\left(\mathrm{ang}\right)}{2}}\]

\[\fitmath{\mathrm{uopt}=\frac{(590)\cdot \sin\left((1.117010721)\right)}{2}\approx 265.1442\,\mathrm{m/s}}\]

\item Se evalúa el máximo ideal.

\[\fitmath{\mathrm{wmax}=\frac{{\mathrm{Vel}}^{2}\cdot {\sin\left(\mathrm{ang}\right)}^{2}}{2}}\]

\[\fitmath{\mathrm{wmax}=\frac{{(590)}^{2}\cdot {\sin\left((1.117010721)\right)}^{2}}{2}\approx 140602.9\,\mathrm{J/kg}}\]

\end{enumerate}
